%% ============================================================
%%  variant-onepage.tex — 极简一页简历（应届 / 校招首选）
%%  - 强制 compact 模式 + nofooter，最大化内容密度
%%  - 精选最核心章节：教育 + 技能 + 经历 + 项目
%%  - 去掉摘要 / 荣誉 / 自我评价等非必要章节
%%  - 适合 ATS 解析：简洁结构、最少图形元素
%% ============================================================
\documentclass[zh,compact,nofooter]{resume-pro}

%% ============================================================
%%  config.tex — 简历全局配置文件
%% ============================================================
%%
%%  ⚠️  发布前请务必替换以下占位数据为你自己的真实信息！
%%      - 姓名、职位、关键词
%%      - 电话、邮箱、社交链接
%%      - 头像照片
%%      - Logo（将 logos/ 下的占位图替换为真实 Logo）
%% ============================================================

%% ---------- 1. 主题色 ----------
% 内置：techblue / businessblack / academicgray / energyorange /
%       innovationgreen / professionalpurple / crimson / ocean
\settheme{techblue}
% \customtheme{0F172A}{2563EB}{E6EEFB}   % 自定义：主色 / 次色 / 浅底（HEX）

%% ---------- 2. 个人信息 ----------
%% ⚠️  以下为示例数据，请替换为你自己的真实信息！
\name{Hertz}                           % ← 替换为你的姓名
\tagline{大模型算法工程师 · 模型压缩 · 多模态}
\keywords{LLM \textbullet{} PyTorch \textbullet{} DeepSpeed \textbullet{} RAG \textbullet{} Agent \textbullet{} 模型压缩}

\phone{+86 138-0000-0000}              % ← 替换为你的真实电话
\email{adong@example.com}              % ← 替换为你的真实邮箱
\github{hertz}
\linkedin{hertz}
\website{https://hertz.dev}
\blog{https://blog.hertz.dev}
\location{北京·海淀}
% \birthday{2000-01}
% \wechat{hertz}
% \zhihu{hertz}
% \scholar{https://scholar.google.com/citations?user=XXXX}
% \orcid{0000-0000-0000-0000}

%% ---------- 3. 个人头像 / 品牌 Logo（二选一） ----------
% 圆形 / 方形 / 圆角，默认 circular
\photoshape{circular}
\photo{logos/avatar.png}        % ← 替换为你的真实头像（构建脚本自动从 ../hertz.jpg 生成）
% \headerlogo{logos/brand.png}  % 顶部右侧品牌 Logo（与 \photo 二选一）

%% ---------- 4. 学校 / 公司 Logo 短宏 ----------
% logos/ 目录下提供了 50+ 占位 Logo（带品牌色 + 首字母）；
% 把同名 PNG 替换为真实 Logo 即可无缝升级。
% 不需要的字段保留为空 \def\xxxLogo{} 即可。
% 字体回退链已内置在 resume-pro.cls 中（macOS / Linux / Windows 自动适配）。
\def\tsinghuaLogo{logos/tsinghua.png}
\def\pekingLogo{logos/peking.png}
\def\sjtuLogo{logos/sjtu.png}
\def\zjuLogo{logos/zju.png}
\def\fudanLogo{logos/fudan.png}
\def\ustcLogo{logos/ustc.png}
\def\stanfordLogo{logos/stanford.png}
\def\mitLogo{logos/mit.png}
\def\ethLogo{logos/eth.png}
% 更多中国高校
\def\rucLogo{logos/ruc.png}
\def\cugLogo{logos/cug.png}
\def\swjtuLogo{logos/swjtu.png}
\def\scuLogo{logos/scu.png}
\def\nankaiLogo{logos/nankai.png}
\def\tjuLogo{logos/tju.png}
\def\dlutLogo{logos/dlut.png}
\def\uestcLogo{logos/uestc.png}
\def\sudaLogo{logos/suda.png}
\def\xmuLogo{logos/xmu.png}
\def\hnuLogo{logos/hnu.png}
\def\csuLogo{logos/csu.png}
% 更多国际高校
\def\caltechLogo{logos/caltech.png}
\def\columbiaLogo{logos/columbia.png}
\def\upennLogo{logos/upenn.png}
\def\nyuLogo{logos/nyu.png}
\def\uscLogo{logos/usc.png}
\def\uclaLogo{logos/ucla.png}
\def\uiucLogo{logos/uiuc.png}
\def\gatechLogo{logos/gatech.png}
\def\umichLogo{logos/umich.png}
\def\uwLogo{logos/uw.png}
\def\ucsdLogo{logos/ucsd.png}

\def\tencentLogo{logos/tencent.png}
\def\bytedanceLogo{logos/bytedance.png}
\def\alibabaLogo{logos/alibaba.png}
\def\baiduLogo{logos/baidu.png}
\def\meituanLogo{logos/meituan.png}
\def\huaweiLogo{logos/huawei.png}
\def\xiaomiLogo{logos/xiaomi.png}
\def\didiLogo{logos/didi.png}
\def\nvidiaLogo{logos/nvidia.png}
\def\openaiLogo{logos/openai.png}
\def\googleLogo{logos/google.png}
\def\microsoftLogo{logos/microsoft.png}
\def\amazonLogo{logos/amazon.png}
\def\appleLogo{logos/apple.png}
\def\metaLogo{logos/meta.png}
\def\anthropicLogo{logos/anthropic.png}
\def\deepseekLogo{logos/deepseek.png}
\def\moonshotLogo{logos/moonshot.png}
\def\minimaxLogo{logos/minimax.png}
\def\alibabaCloudLogo{logos/alibaba-cloud.png}

%% ---------- 5. 页脚（fileinfo + version）----------
% 出现在每页底部：左 fileinfo · 中 页码 · 右 version (+ 可选 QR)
% 用 \rpName 引用姓名（在 \name{...} 后自动可用，由 cls 在 \AtBeginDocument 时缓存）
\fileinfo{\faCopyright\ 2025--2026 \rpName\ \textbullet\ Built with \LaTeX}
\version{v3.9.0 \textbullet\ \today}

%% ---------- 6. 页脚 QR 码（可选）----------
% 设置后会在页脚右侧 version 前嵌入一个小 QR 码，指向你设定的 URL。
% 适合面试时扫码直达个人主页 / GitHub / 作品集。
% \qrlink{https://hertz.dev}

%% ---------- 7. 最大页数限制（可选）----------
% 设置后，编译结束时若超出页数会在终端输出警告。
% \setMaxPages{1}   % 应届一页简历建议设为 1


\settheme{techblue}
\version{One-Page \textbullet\ \today}
\setMaxPages{1}

\begin{document}

\makeheader

%% ===== 专业技能（精简为一行标签）=====
\sectionTitle{专业技能}{\faLaptopCode}

\skillrow{编程}{Python \textbullet{} C++ / CUDA \textbullet{} Bash}
\skillrow{大模型}{\skilltagAccent{LLM} \skilltagAccent{SFT} \skilltagAccent{RLHF/DPO} \skilltag{LoRA} \skilltag{RAG} \skilltag{Agent} \skilltag{Quantization}}
\skillrow{框架}{PyTorch \textbullet{} DeepSpeed \textbullet{} vLLM \textbullet{} LangChain \textbullet{} Docker / K8s}

%% ===== 教育背景 =====
\sectionTitle{教育背景}{\faGraduationCap}

\educationItem[\tsinghuaLogo]{清华大学 \textbullet{} 计算机系}{2023.09 -- 2026.06}{计算机科学与技术 \textbullet{} 硕士}{排名前 5\% \textbullet{} 国家奖学金}
\educationItem[\tsinghuaLogo]{清华大学 \textbullet{} 计算机系}{2019.09 -- 2023.06}{计算机科学与技术 \textbullet{} 学士}{GPA 3.85/4.00 \textbullet{} ETH 交换}

%% ===== 实习经历 =====
\sectionTitle{实习经历}{\faBriefcase}

\experienceItem[\bytedanceLogo]{字节跳动 \textbullet{} 豆包大模型团队}{2024.06 -- 2024.12}{大模型算法实习生}{%
\begin{itemize}
  \item CoT 数据构造 5B tokens 持续预训练，HumanEval/MATH 平均 \textbf{+5--8 pp}
  \item 多任务 SFT + DPO（LLaMA-3-8B），满意度 \textbf{87\%}，benchmark \textbf{+3 pp}
  \item ReAct + ToT 混合推理框架，工具调用 \textbf{72\% → 89\%}
\end{itemize}
}

\experienceItem[\tencentLogo]{腾讯 \textbullet{} 微信 AI 团队}{2023.06 -- 2023.12}{算法工程实习生}{%
\begin{itemize}
  \item 推理服务优化，P99 \textbf{3.2s → 1.8s}，QPS \textbf{+65\%}
  \item INT8 量化方案，性能保留 \textbf{98\%}，显存 \textbf{−55\%}
\end{itemize}
}

%% ===== 项目精选 =====
\sectionTitle{项目经历}{\faProjectDiagram}

\projectCard{自我纠错型 RAG 系统}{2024.01 -- 2024.05}
{企业知识库 Agent，召回率 65\% → 85\%}{%
\begin{itemize}
  \item 迭代式检索 + 混合检索 + 交叉编码器重排序，MS MARCO \textbf{+15\%}
\end{itemize}
}{\faGithub\ github.com/hertz/self-rag-agent}

\projectCard{LLM 推理加速框架 \textbullet{} FastInfer}{2023.03 -- 2023.06}
{基于 vLLM 二开的低延迟推理引擎}{%
\begin{itemize}
  \item PagedAttention + KV Cache 复用，A100 吞吐 \textbf{2400 tok/s}，提升 \textbf{1.8×}
\end{itemize}
}{\faGithub\ github.com/hertz/fastinfer}

%% ===== 代表性成果（一行式）=====
\sectionTitle{代表性成果}{\faTrophy}

\inlineEntry{\faFile*}{NeurIPS 2024 一作（CCF-A）}{层级知识蒸馏}
\inlineEntry{\faTrophy}{Kaggle LLM 金牌 Top 1\%}{2023}
\inlineEntry{\faGithub}{开源 5000+ Stars}{Agent-RL-Framework}

\end{document}
